% Hướng dẫn sử dụng cho tổ sản xuất (tổ trưởng, thợ) và KCS, ERP Sao Việt Nhật.
% Biên dịch: xelatex huong-dan-viec-cua-to-va-kcs.tex (chạy hai lần cho mục lục).
% Người đọc: docs/personas/to-san-xuat.md, docs/personas/kcs.md. Khai báo nằm ở
% ../huong-dan-cau-hinh/cau-hinh-he-thong.tex, phần lên lệnh ở ../huong-dan-ke-hoach-sx/huong-dan-ke-hoach-san-xuat.tex.
% Ảnh chụp bằng kịch bản chạy lại được (kich-ban/chup-ui.mjs + c02..c09), dữ liệu dựng bằng thao tác thật trên dev.
\documentclass[11pt,a4paper]{article}
\usepackage{fontspec}
\usepackage[vietnamese]{babel}
\setmainfont{Times New Roman}
\newfontfamily\tieude{Arial}
\usepackage[a4paper,margin=1.9cm,top=1.7cm]{geometry}
\usepackage{graphicx}
\usepackage{enumitem}
\usepackage{xcolor}
\usepackage{titlesec}
\usepackage{tocloft}
\usepackage{needspace}
\usepackage{fancyhdr}
\usepackage[hidelinks,destlabel=true]{hyperref}
\usepackage[most]{tcolorbox}
\usepackage{booktabs}
\usepackage{tabularx}
\graphicspath{{hinh/}}

\definecolor{gach}{HTML}{C2410C}
\definecolor{xamnen}{HTML}{F4F4F5}
\definecolor{xamchu}{HTML}{52525B}

\setlist{nosep,leftmargin=1.5em}
\setlength{\parskip}{0.4em}
\setlength{\parindent}{0pt}
\setlength{\emergencystretch}{2em}

\titleformat{\section}{\tieude\Large}{Chương \thesection.}{0.4em}{}[{\color{black!25}\titlerule[0.4pt]}]
\titlespacing*{\section}{0pt}{1.2em}{0.6em}
\titleformat{\subsection}{\tieude\large}{\thesubsection}{0.6em}{}
\titlespacing*{\subsection}{0pt}{1.1em}{0.3em}
\renewcommand{\cftsecpresnum}{Chương }
\setlength{\cftsecnumwidth}{5.9em}
\renewcommand{\cftsecfont}{\tieude}
\renewcommand{\cftsecpagefont}{\tieude}
\setlength{\cftbeforesecskip}{0.5em}
\addto\captionsvietnamese{\renewcommand{\contentsname}{Mục lục}}
\renewcommand{\cfttoctitlefont}{\tieude\large}
\setlength{\cftbeforetoctitleskip}{0pt}
\setlength{\cftaftertoctitleskip}{0pt}

\pagestyle{fancy}
\fancyhf{}
\fancyhead[L]{\tieude\footnotesize\color{xamchu}Hướng dẫn sử dụng Việc của tổ và KCS}
\fancyhead[R]{\tieude\footnotesize\color{xamchu}ERP Sao Việt Nhật}
\fancyfoot[C]{\tieude\footnotesize\thepage}
\renewcommand{\headrulewidth}{0.4pt}

\newtcolorbox{boicanh}{colback=xamnen,colframe=xamnen,boxrule=0pt,arc=3pt,
  left=9pt,right=9pt,top=6pt,bottom=6pt,before skip=4pt,after skip=8pt}

\definecolor{vang}{HTML}{FEF08A}
\newcommand{\oo}[1]{{\setlength{\fboxsep}{1.2pt}\colorbox{vang}{\strut #1}}}
\newcommand{\duong}[1]{{\tieude\small #1}}
\newcommand{\so}[1]{{\tieude\bfseries\small\textcolor{white}{\colorbox{red!80!black}{\kern-2pt#1\kern-2pt}}}}
\newcommand{\khung}[2]{{\setlength{\fboxsep}{0pt}\setlength{\fboxrule}{0.5pt}%
  \fcolorbox{black!20}{white}{\includegraphics[width=\dimexpr#1-1pt\relax]{#2}}}}
\newcommand{\q}[1]{\textquotedbl#1\textquotedbl{}}
\newcommand{\canh}[3]{%
  \par\noindent
  \begin{minipage}[t]{#1\linewidth}\vspace{0pt}\khung{\linewidth}{#2}\end{minipage}\hfill
  \begin{minipage}[t]{\dimexpr0.97\linewidth-#1\linewidth\relax}\vspace{0pt}\setlength{\parskip}{0.5em}#3\end{minipage}\par\medskip}
\newcommand{\rong}[2][\linewidth]{\par\noindent\begin{center}\khung{#1}{#2}\end{center}\vspace{-0.6em}}
% Ảnh cắt bớt: \cat[rộng]{trái dưới phải trên (px)}{ảnh}.
\newcommand{\cat}[3][\linewidth]{\par\noindent\begin{center}{\setlength{\fboxsep}{0pt}\setlength{\fboxrule}{0.5pt}%
  \fcolorbox{black!20}{white}{\includegraphics[width=\dimexpr#1-1pt\relax,trim=#2,clip]{#3}}}\end{center}\vspace{-0.6em}}
% Tên tài liệu khác.
\newcommand{\tl}[1]{\textit{#1}}
% Liên kết sang một mục của tài liệu Cấu hình và tài liệu Kế hoạch SX (đích là tên \label bên đó, nhờ destlabel).
\newcommand{\cauhinh}{../huong-dan-cau-hinh/cau-hinh-he-thong.pdf}
\newcommand{\kehoach}{../huong-dan-ke-hoach-sx/huong-dan-ke-hoach-san-xuat.pdf}
\definecolor{lienket}{HTML}{1E40AF}
\newcommand{\kbn}[2]{\href{\cauhinh\##1}{\textcolor{lienket}{#2}}}
\newcommand{\kb}[2]{\tl{Hướng dẫn cấu hình hệ thống}, mục \kbn{#1}{#2}}
\newcommand{\khn}[2]{\href{\kehoach\##1}{\textcolor{lienket}{#2}}}
\newcommand{\kh}[2]{\tl{Hướng dẫn sử dụng Kế hoạch sản xuất}, mục \khn{#1}{#2}}

\begin{document}

\thispagestyle{empty}
{\tieude\small\color{xamchu} ERP Sao Việt Nhật}\par\vspace{3pt}
{\tieude\fontsize{22}{26}\selectfont Hướng dẫn sử dụng Việc của tổ và KCS}\par\vspace{4pt}
{\small Nhận việc, bắt đầu, ghi sản lượng, giao tổ sau, kiểm hàng, gửi kho, đóng lệnh}\par\vspace{2pt}
{\small\color{xamchu} Cập nhật 05/11/2026}\par\vspace{6pt}
\noindent{\color{black!25}\rule{\linewidth}{0.4pt}}\par\vspace{4pt}

Tài liệu cho người dưới xưởng: tổ trưởng và thợ của các tổ in, cán, bồi, bế, gỡ, dán, thành phẩm, cắt, và người
làm KCS. Đi theo đúng đường của một việc: lệnh phát hành xuống, tổ nhận việc, bắt đầu, ghi sản lượng, giao tổ sau,
tới bước cuối thì KCS kiểm, gửi thành phẩm vào kho và đóng lệnh. Người đọc rành nghề hơn ai hết nên tài liệu
không giải nghĩa tờ in, con, khuôn bế, chỉ nói phần mềm làm gì.

Mỗi tổ là một dòng trong ma trận quyền, nhóm Tổ sản xuất, có bốn ô:

\begin{itemize}
\item \oo{Xem}: thấy tổ ở thanh bên, mở được màn Việc của tổ.
\item \oo{Thực hiện lệnh}: giao người, bắt đầu, tạm dừng, kết thúc, ghi sản lượng, nhận và trả khuôn, đề xuất giao,
yêu cầu cấp vật tư.
\item \oo{Xác nhận sản lượng}: xác nhận đã nhận hàng của tổ trước, sửa số một lần ghi.
\item \oo{Kho}: yêu cầu nhập kho phần vật tư còn dư.
\end{itemize}

Vai thợ có Xem và Thực hiện lệnh. Vai tổ trưởng có thêm Xác nhận sản lượng. Màn KCS là một dòng riêng với ba ô
\oo{Xem}, \oo{Thao tác}, \oo{Đóng lệnh thiếu}, viết ở chương \ref{kcs}. Thiếu ô nào thì nút tương ứng không hiện
hoặc bấm bị từ chối. Quản trị bật quyền ở \kb{phan-quyen-ma-tran}{Mở ma trận quyền} và gán vai ở mục
\kbn{phan-quyen-gan-vai}{Gán vai và lập nhóm}. Chữ màu xanh là liên kết, bấm vào mở đúng mục bên tài liệu đó.

Ảnh chụp từ phần mềm thật, dữ liệu là lệnh LSX26-0007 (Folder A4 cán mờ ép kim, 1.000 cái, khách Café Katinat)
chạy đủ bảy bước từ In tới đóng lệnh. Ảnh chụp bằng tài khoản quản trị nên đầu màn ghi tên Admin và có thể hiện
nhiều nút hơn màn của một thợ. Ô đỏ đánh số là chỗ cần bấm hoặc cần đọc, số khớp với lời giải thích bên cạnh.

\bigskip
\tableofcontents

% ============================================================
\Needspace*{27\baselineskip}
\section{Đường đi của việc dưới xưởng}\label{duong-di}

Mỗi hàng của sơ đồ là một bên làm việc, mỗi ô là một bước. Chữ trên mũi tên là điều kiện hoặc thứ được chuyển
sang bước sau.

\cat[0.95\linewidth]{0 330 0 90}{vtt-00-luong.png}

Kế hoạch sản xuất phát hành lệnh (xem \kh{phat-hanh}{Phát hành}). Lúc đó mỗi bước của lệnh thành một việc ở
đúng tổ được giao, kho nhận yêu cầu xuất vật tư. Tổ trưởng giao người, nhận khuôn nếu bước cần khuôn, rồi bấm
\oo{Bắt đầu}. Thợ ghi sản lượng từng lần, ghi xong phần mềm tự đề xuất giao lần đó sang tổ sau. Tổ sau bấm
\oo{Xác nhận} là đã nhận hàng và bắt đầu được bước của mình. Cứ thế tới bước cuối.

Ở bước cuối, KCS kiểm hàng: bao nhiêu đạt, bao nhiêu lỗi, lỗi do công đoạn nào gây ra. Hàng đạt gửi vào kho thành
phẩm, kho ghi sổ nhập, KCS đóng lệnh. Giữa chừng KCS thấy lỗi ở bước nào thì ghi lỗi ngay ở bước đó, tổ gây lỗi
thấy liền ở tab KCS của việc.

Bước thuê ngoài không thuộc tổ nào nên không hiện ở Việc của tổ. Kế hoạch sản xuất mang hàng đi và nhận về (xem
\kh{mot-phan}{Thuê ngoài một phần}), tổ của bước sau chỉ việc xác nhận nhận hàng như nhận từ tổ trong xưởng.

% ============================================================
\Needspace*{30\baselineskip}
\section{Màn Việc của tổ}\label{man-viec}

Mỗi tổ có một màn riêng. Mở bằng tên tổ ở nhóm \duong{Tổ sản xuất} của thanh bên trái. Người có quyền ở nhiều tổ
thấy nhiều tên, mỗi tên một màn.

\rong{vtt-10-man-viec-cua-to.png}

\so{1} Tên tổ ở thanh bên. Tổ con nằm dưới tổ cha, tổ cha chỉ là tiêu đề nhóm.
\so{2} Ba kiểu xem: \oo{Bảng}, \oo{Lịch}, \oo{Sản lượng}. Nút \oo{Báo dừng máy} bên cạnh viết ở mục \ref{bao-dung}.
\so{3} Ô tìm theo mã lệnh, khách, công đoạn, máy. Cạnh đó là chọn kỳ, \oo{Lọc} (lọc thêm theo trạng thái, lệnh gấp,
hàng đang tới đâu, lý do chưa bắt đầu được, máy, khách, người làm), \oo{Sắp xếp} và \oo{Cột}.
\so{4} Chip trạng thái kèm số việc: Tất cả, Đang chạy, Tạm dừng, Chờ làm, Hoàn thành. Có hàng tổ trước giao chưa
xác nhận thì hiện thêm chip hàng chờ nhận, có lỗi KCS tổ chưa xem thì hiện chip \oo{Lỗi KCS chờ xem} (mục
\ref{to-xem-loi}).
\so{5} Cột Vật tư: \oo{Sẵn sàng} là kho đã cấp đủ, còn lại là đang chờ kho hoặc kho mới cấp một phần.
\so{6} Cột Tốt / cần: số tốt đã ghi trên số phải làm, thanh xanh là phần đã xong.
\so{7} Cột Hàng đang tới đâu: dải chấm là các bước của lệnh, chấm to là tổ mình, dòng chữ dưới nói hàng đang ở đâu.

Mỗi dòng là một việc. Cột Công đoạn có chip \oo{Máy} hay \oo{Tổ} (việc chạy máy hay làm tay) và chip khuôn nếu bước
cần khuôn. Cột Trạng thái ghi giờ xong hoặc tên người đang làm. Bấm vào dòng để mở ngăn của việc.

\Needspace*{26\baselineskip}
\subsection{Kiểu Lịch và kiểu Sản lượng}\label{kieu-xem}

\rong{vtt-11-kieu-lich.png}

Kiểu \oo{Lịch} xếp việc theo ngày: Tuần, 2 tuần hoặc Tháng, nút Hôm nay và mũi tên để lùi tới. Ô nét đứt là giờ
kế hoạch xếp, ô đặc là phiên tổ đã chạy thật. Dòng đầu đếm số việc, giờ tổ đã chạy so với giờ kế hoạch và ca làm
việc. Cột \oo{Hàng chờ} bên phải giữ những việc chưa định giờ hoặc nằm ngoài khoảng ngày đang xem.

\rong{vtt-12-kieu-san-luong.png}

Kiểu \oo{Sản lượng} là sổ sản lượng của tổ: mỗi lệnh một khối, có giấy và khổ tờ, rồi từng công đoạn tổ đã làm,
lần ghi, việc phát sinh và số lượng. Chọn kỳ theo ngày bắt đầu, mặc định là Tháng này.

\Needspace*{29\baselineskip}
\subsection{Ngăn của một việc}\label{ngan-viec}

\rong[0.85\linewidth]{vtt-13-ngan-viec.png}

Đầu ngăn là mã lệnh, tên sản phẩm, khách, tên công đoạn và trạng thái. Nút việc (Bắt đầu, Tạm dừng, Kết thúc, dấu
ba chấm) nằm góc phải, đổi theo trạng thái.

\so{1} Sáu tab: \oo{Quy cách} (giấy, khổ, màu, số kẽm, kế hoạch khối lượng và giờ), \oo{Sản lượng} (chương
\ref{ghi-san-luong}), \oo{Giao và nhận} (chương \ref{giao-nhan}), \oo{Vật tư} (chương \ref{vat-tu}), \oo{KCS}
(mục \ref{to-xem-loi}), \oo{Lịch sử} (ai bấm gì, lúc nào).

\so{2} \oo{Hàng đang tới đâu}: các bước của lệnh từ trên xuống, chữ \q{tổ mình} đánh dấu bước của tổ, cuối dải là
Nhập kho thành phẩm. Câu đầu nói gọn hàng đang ở đâu.

\so{3} \oo{Tiến độ}: số tốt trên số cần, Nhận vào, Lỗi, Đã giao đi, Chờ giao, Còn phải làm. KCS bắt lỗi ở bước
này thì có thêm dòng KCS bắt lỗi.

\so{4} \oo{Việc}: máy (kèm phiếu sửa chữa đang mở của máy nếu có), thợ được giao, khuôn, dặn dò của kế hoạch.
Dưới cùng là \oo{Tệp của lệnh}, file kế hoạch đính kèm.

Việc chưa bắt đầu được thì đầu cột phải có thêm khung vàng \oo{Chưa bắt đầu được vì} (mục \ref{chua-bat-dau}).

\Needspace*{32\baselineskip}
\subsection{Trên điện thoại}\label{dien-thoai}

\par\noindent\begin{center}\khung{0.34\linewidth}{vtt-14-dt-danh-sach.png}\hspace{2em}\khung{0.34\linewidth}{vtt-15-dt-ngan-viec.png}\end{center}\vspace{-0.6em}

Màn hẹp thì mỗi việc thành một thẻ: công đoạn, trạng thái, vật tư, tốt / cần, hàng đang tới đâu, máy, quy cách,
định mức vật tư. Bộ lọc và chip trạng thái vẫn ở đầu màn, kéo ngang để xem hết chip.

Bấm thẻ để mở việc. Ngăn mở hết màn, có thêm tab \oo{Tổng quan} gom hàng đang tới đâu, tiến độ và phần việc mà
máy tính để ở cột phải. Mũi tên góc trái để quay về danh sách.

Thanh dưới cùng là các màn của riêng người dùng (Của tôi, Công của tôi, Nghỉ phép, Phiếu lương) và \oo{Menu} để
mở thanh bên.


% ============================================================
\Needspace*{30\baselineskip}
\section{Giao người và bắt đầu}\label{giao-nguoi}

Lệnh vừa phát hành thì cả tổ tự được giao vào việc. Tổ trưởng chỉnh lại cho đúng người đứng việc này.

\rong[0.85\linewidth]{vtt-21-giao-nguoi.png}

\so{1} \oo{Bắt đầu}: bấm là việc chuyển sang Đang chạy, giờ máy bắt đầu tính. Không có hộp hỏi lại.
\so{2} Mỗi thợ là một chip. Bấm dấu x trên chip để rút người ra.
\so{3} \oo{+ Giao người} mở ô tìm theo tên hoặc mã. Cả tổ đã được giao thì ô báo \q{Không còn ai để giao.}

Giao người và rút người cần ô Thực hiện lệnh. Người được giao là người hiện sẵn ở ô \q{Người làm lần này} khi ghi
sản lượng (chương \ref{ghi-san-luong}).

\Needspace*{29\baselineskip}
\subsection{Vì sao chưa bắt đầu được}\label{chua-bat-dau}

\rong[0.85\linewidth]{vtt-24-chua-bat-dau-duoc.png}

Nút Bắt đầu mờ thì khung vàng \oo{Chưa bắt đầu được vì} nói lý do. Có bốn lý do:

\begin{itemize}
\item Chưa giao thợ: việc chưa có ai. Tổ trưởng giao người.
\item Chưa nhận khuôn: bước cần khuôn mà tổ chưa bấm Đã nhận khuôn (mục \ref{khuon}).
\item Chưa nhận hàng bước trước: tổ trước chưa giao, hoặc đã giao mà tổ mình chưa bấm Xác nhận ở tab Giao và nhận.
Ảnh trên là Ép kim đang chờ hàng cán màng thuê ngoài về.
\item Chưa sẵn sàng vật tư: kho chưa cấp, hoặc mới cấp một phần (chương \ref{vat-tu}).
\end{itemize}

Bộ lọc \oo{Chưa bắt đầu được vì} ở đầu màn lọc đúng theo bốn lý do này, để tổ trưởng thấy ngay việc nào đang kẹt
ở đâu.

\Needspace*{29\baselineskip}
\subsection{Nhận khuôn và trả khuôn}\label{khuon}

\rong[0.85\linewidth]{vtt-25-nhan-khuon.png}

Bước cần khuôn có dòng \oo{Khuôn} ở cột phải: mã khuôn, tên khuôn và tình trạng của khuôn. Lấy khuôn ra máy rồi
bấm \so{1} \oo{Đã nhận khuôn}. Chip khuôn đổi thành \q{đã nhận}, Bắt đầu hết bị chặn vì khuôn.

\rong[0.85\linewidth]{vtt-26-tra-khuon.png}

Làm xong việc, nút \so{1} \oo{Trả khuôn về kệ} mới hiện. Cất khuôn rồi bấm để khuôn về lại kệ, việc khác dùng
khuôn này mới nhận được.

\Needspace*{14\baselineskip}
\subsection{Đang chạy}\label{dang-chay}

\rong{vtt-23-dang-chay.png}

Bắt đầu xong, \so{1} chip trạng thái ghi \q{Đang chạy từ} và giờ bấm. Góc phải đổi thành \so{2} \oo{Tạm dừng} và
\so{3} \oo{Kết thúc}, viết ở chương \ref{dung-ket-thuc}.

% ============================================================
\Needspace*{34\baselineskip}
\section{Ghi sản lượng}\label{ghi-san-luong}

Làm được bao nhiêu thì ghi bấy nhiêu, ghi nhiều lần trong một việc. Mỗi lần ghi là căn cứ tính lương khoán cho
người làm lần đó. Mở tab \oo{Sản lượng}, bấm \oo{+ Ghi sản lượng}.

\canh{0.46}{vtt-30-ghi-san-luong.png}{%
\so{1} \oo{Số lượng làm được}, theo đơn vị ra của bước (tờ in, con, cái). Góc phải ghi \q{còn ghi được} bao nhiêu.
Nút \oo{Tối đa} điền đúng số còn lại.

\so{2} Nút cộng nhanh \oo{+10}, \oo{+50}, \oo{+100}, \oo{+1.000} và \oo{Về 0}. Thanh dưới cho biết sau lần này
được bao nhiêu trên số kế hoạch.

\so{3} \oo{Người làm lần này}: người đang được giao, tick sẵn cả. Bỏ tick người không đứng ca này. Người đang làm
việc khác cùng lúc thì có chữ nhắc bên cạnh tên.

\so{4} \oo{Thêm người tổ khác}: mượn người tổ khác vào làm lần này.

\so{5} \oo{Ghi sản lượng}, hoặc bấm Ctrl + Enter.}

\canh{0.46}{vtt-30b-ghi-san-luong-duoi.png}{%
Kéo xuống dưới hộp:

\so{6} \oo{Khoán}: công thức tính tiền của bước, và \oo{Bậc giá} lệnh này rơi vào (bước in theo số màu). Công thức
khai ở \kb{danh-muc-khoan}{Tiền khoán của công đoạn}.

\so{7} \oo{Tuỳ chọn lần này}: tick thì giá cả lần ghi đổi theo, ví dụ in mực UV.

\so{8} \oo{Việc phát sinh}: thay kẽm, xuống bài do lỗi máy. Tính tiền riêng, không cộng vào sản lượng. Có làm thì
gõ số, không làm thì để trống.

\so{9} Dòng cuối báo ghi xong sẽ tự đề xuất giao lần này sang bước nào. Bước sau thuộc tổ khác thì lần này nằm ở
Giao đi chờ tổ đó xác nhận (chương \ref{giao-nhan}).}

Mục Ghi chú để ghi thêm điều cần nhớ của lần này. Số ghi không vượt được số còn phải làm. KCS bắt lỗi ở bước này
thì số còn ghi được trừ luôn phần lỗi.

\Needspace*{29\baselineskip}
\subsection{Bảng các lần ghi}\label{cac-lan-ghi}

\rong[0.85\linewidth]{vtt-31b-mo-lan.png}

Mỗi lần ghi là một dòng \q{Lần N}: lúc ghi, người làm, máy, số tốt, số hỏng. Bấm một dòng để mở ra: đơn giá khoán
và bậc giá của lần đó, cùng nút \oo{Sửa số}.

\Needspace*{22\baselineskip}
\subsection{Sửa số đã ghi}\label{sua-so}

\rong[0.75\linewidth]{vtt-32-sua-so.png}

Ghi nhầm hay đếm lại thấy khác thì tổ trưởng (ô Xác nhận sản lượng) bấm \oo{Sửa số} ở dòng lần đó. Hộp có \oo{Số
đang ghi} và \oo{Số sau khi sửa}, nút trừ cộng 100, 10, 1, người làm, tuỳ chọn, và \oo{Lý do sửa} (chọn nhanh Đếm
lại, Gõ nhầm). Lần giao đi từ lần ghi này đổi theo số mới. Số sửa và lý do vào tab Lịch sử của việc.

% ============================================================
\Needspace*{31\baselineskip}
\section{Giao và nhận giữa các tổ}\label{giao-nhan}

Tab \oo{Giao và nhận} có hai phần: \oo{Nhận vào} (hàng tổ trước giao cho mình) và \oo{Giao đi} (hàng mình giao cho
bước sau).

\rong[0.85\linewidth]{vtt-40-giao-di.png}

Ảnh là bước In sau hai lần ghi 150 và 100 tờ. Bước đầu tiên không nhận hàng của ai, Nhận vào ghi \q{từ kho}: vật tư
lấy theo phiếu xuất ở tab Vật tư. Giao đi có hai dòng, mỗi lần ghi một dòng, đang \oo{chờ xác nhận} của bên nhận.
Bấm \oo{Chọn lại} ở một dòng để đổi bước nhận (khi lệnh có nhiều nhánh). Nút \oo{Đề xuất giao} để giao tay phần
hàng tốt chưa giao, dùng khi ghi sản lượng không tự đề xuất được, ví dụ bước sau có nhiều nhánh.

\Needspace*{29\baselineskip}
\subsection{Xác nhận nhận hàng}\label{xac-nhan}

\rong[0.85\linewidth]{vtt-41-nhan-vao.png}

Tổ sau thấy chấm đỏ trên tab Giao và nhận. Đếm hàng thực nhận rồi tổ trưởng bấm \so{1} \oo{Xác nhận}. Bấm là xong,
không có hộp hỏi lại. Trạng thái dòng đổi thành \oo{đã xác nhận}, khung \q{Chưa bắt đầu được vì} mất lý do chờ hàng.

Hai bước liền nhau cùng một tổ thì không phải xác nhận: dòng nhận ghi \q{Tự nhận (cùng tổ)} và đã xác nhận sẵn.

\rong[0.85\linewidth]{vtt-43-da-nhan.png}

Hàng thuê ngoài về (ảnh trên là 250 tờ cán màng về bước Ép kim) cũng xác nhận y như vậy. Cột Từ ghi tên nhà gia
công.

\Needspace*{10\baselineskip}
\subsection{Bước cuối}\label{buoc-cuoi}

Bước cuối của lệnh không giao cho tổ nào. Phần Giao đi ghi \q{Bước cuối, hàng tốt} chờ KCS kiểm. KCS kiểm xong
mới gửi vào kho (chương \ref{kcs}).

% ============================================================
\Needspace*{31\baselineskip}
\section{Vật tư của việc}\label{vat-tu}

Lúc phát hành, phần mềm tự lập yêu cầu xuất kho cho vật tư của từng bước. Tab \oo{Vật tư} cho tổ xem kho đã cấp
tới đâu và xin thêm khi thiếu.

\rong[0.85\linewidth]{vtt-50-tab-vat-tu.png}

Bảng \oo{Vật tư theo kế hoạch}: mỗi món một dòng với Kế hoạch, Đã xin, Kho đã xuất, Nhập lại, Thực dùng và Tình
trạng (\oo{Khớp} là kho đã xuất đủ số đã xin). Bấm mũi tên đầu dòng để xem chi tiết từng lần xuất. Bảng \oo{Đề nghị
và phiếu kho}: các lần xin, chiều (xuất hay nhập), cần lúc, người gửi, trạng thái, phiếu kho đã lập. Dòng \q{Tự lập
khi phát hành} là yêu cầu phần mềm lập sẵn.

Góc phải có một nút xin, đổi theo tình hình:

\begin{itemize}
\item \oo{Yêu cầu cấp vật tư}: bước chưa có lần xin nào.
\item \oo{Sửa đề nghị}: lần xin gần nhất kho chưa xử lý, còn sửa được.
\item \so{2} \oo{Yêu cầu bổ sung}: đã xin đủ kế hoạch mà vẫn thiếu, ví dụ hư giấy nhiều hơn bù hao.
\end{itemize}

Và \so{1} \oo{Yêu cầu nhập kho}: trả phần vật tư còn dư về kho, cần ô Kho của tổ.

\canh{0.5}{vtt-51-bo-sung.png}{%
Hộp \oo{Yêu cầu bổ sung} liệt kê các món của bước với số đã xin và tồn kho. Gõ \oo{Số xin thêm} cho món đang
thiếu, món không thiếu để trống. Chọn \oo{Cần xuất lúc} (Hôm nay, Mai, +3 ngày, Lúc bắt đầu việc), ghi chú nếu cần,
rồi bấm \oo{Gửi kho}. Lệnh, công đoạn, tổ khoá theo việc, không phải chọn.

Hộp \oo{Yêu cầu nhập kho} giống vậy, cột là Kho đã xuất và \oo{Số lượng nhập}: điền số còn thừa trả về kho, món
không trả để 0.}

Yêu cầu gửi đi nằm ở màn Yêu cầu nhập xuất của kho. Kho lập phiếu, ghi sổ thì cột Kho đã xuất và Tình trạng ở tab
này đổi theo. Vật tư chưa đủ thì Bắt đầu bị chặn (mục \ref{chua-bat-dau}).

% ============================================================
\Needspace*{17\baselineskip}
\section{Tạm dừng, kết thúc, đổi máy, báo dừng máy}\label{dung-ket-thuc}

\canh{0.45}{vtt-60-tam-dung.png}{%
\oo{Tạm dừng} mở hộp hỏi \oo{Vì sao dừng}, bắt buộc ghi, ví dụ thay dao cắt khoảng 20 phút. Phiên chạy đóng lại,
giờ máy ngừng tính từ lúc bấm. Việc về trạng thái Tạm dừng, nút góc phải đổi thành \oo{Tiếp tục}. Bấm Tiếp tục là
mở phiên chạy mới.}

\Needspace*{14\baselineskip}
\subsection{Kết thúc}\label{ket-thuc}

\rong{vtt-61b-hoan-thanh.png}

\oo{Kết thúc} đóng việc, chip đổi thành \q{Hoàn thành} kèm giờ. Bấm là xong, không hỏi lại. Kết thúc được khi mọi
hàng tốt đã ghi đều đã đề xuất giao sang bước sau. Còn hàng chưa giao thì phần mềm báo, tổ vào tab Giao và nhận bấm
Đề xuất giao trước. Việc có khuôn thì sau Kết thúc mới trả khuôn (mục \ref{khuon}).

\Needspace*{9\baselineskip}
\subsection{Đổi máy}\label{doi-may}

\rong[0.6\linewidth]{vtt-62-menu-doi-may.png}

\canh{0.45}{vtt-62b-doi-may.png}{%
Nút ba chấm cạnh Kết thúc có \oo{Đổi máy}: chuyển việc sang máy khác cùng chạy được công đoạn này. Hộp ghi \oo{Máy
đang dùng}, ô \oo{Đổi sang máy} và \oo{Lý do đổi} (không bắt buộc). Tổ không còn máy nào khác chạy công đoạn này thì
ô chọn báo \q{Không có máy đổi được}, việc chờ máy sửa xong.}

\Needspace*{28\baselineskip}
\subsection{Báo dừng máy}\label{bao-dung}

\canh{0.45}{vtt-63-bao-dung-may.png}{%
Nút \oo{Báo dừng máy} ở đầu màn Việc của tổ. Hộp gồm:

\begin{itemize}
\item \oo{Ngày}, mặc định hôm nay.
\item \oo{Lý do}: \oo{Máy hư} hoặc \oo{Hết hàng}.
\item Máy hư thì chọn máy. Dưới tên máy ghi việc đang chạy trên máy.
\item \oo{Máy đang dừng, không chạy được}: tick thì việc đang chạy trên máy tự tạm dừng, giờ máy ngừng tính từ
lúc gửi.
\item \oo{Mức độ theo bạn thấy}: Nhẹ, Trung bình, Nghiêm trọng. Tổ sửa chữa đánh giá lại khi nhận.
\item \oo{Triệu chứng}: kể đúng cái mình thấy.
\item \oo{Tính ngày dừng cho lương}: tick thì ngày đó tổ ăn trọn chỉ tiêu ngày, cộng thêm tiền các lần ghi trong
ngày nếu có.
\end{itemize}

Bấm \oo{Gửi tổ sửa chữa}. Phiếu sửa chữa hiện ngay ở dòng Máy trong cột phải của các việc chạy máy đó.}

Gửi phiếu sửa chữa cần ô Yêu cầu của màn Kỹ thuật máy. Thợ chưa có ô này thì nhờ tổ trưởng gửi.

% ============================================================
\Needspace*{26\baselineskip}
\section{KCS kiểm hàng}\label{kcs}

Màn \duong{Sản xuất > KCS} là màn của người làm KCS, không theo tổ: thấy lệnh của mọi tổ. Ô quyền ở dòng KCS của
ma trận (xem \kb{danh-muc-kcs-quyen}{Ai được kiểm KCS}):

\begin{itemize}
\item \oo{Xem}: mở màn, xem các lần kiểm và báo cáo lỗi, xuất Excel.
\item \oo{Thao tác}: kiểm hàng, ghi lỗi, gửi nhập kho, đóng lệnh khi đã đủ.
\item \oo{Đóng lệnh thiếu}: đóng khi còn thiếu số hay còn việc dở, mở lại lệnh đã đóng.
\end{itemize}

Tiêu chí kiểm của từng công đoạn khai ở \kb{danh-muc-kcs}{Tiêu chí KCS}. Lệnh chép bộ tiêu chí lúc phát hành.

\rong{vtt-70-kcs-danh-sach.png}

\so{1} Chip lọc: Tất cả, \oo{Chờ kiểm} (bước cuối có hàng tổ đã làm mà chưa kiểm), \oo{Chưa gửi kho} (đã kiểm đạt
mà chưa gửi kho), \oo{Chờ đóng lệnh}, \oo{Đã đóng}.
\so{2} Ô tìm theo mã lệnh, sản phẩm, khách. Cạnh đó là kỳ, Lọc, Cột.
\so{3} Mỗi dòng là một lệnh: Chuỗi là dãy ô màu của các bước (xanh lá là xong, xanh dương là đang làm, xám là chưa
làm, đỏ là gây lỗi), \oo{Đã kiểm đạt / cần}, số \oo{Chờ kiểm}, và \oo{Lỗi do lệnh này gây} theo công đoạn.

\Needspace*{33\baselineskip}
\subsection{Sơ đồ bước của lệnh}\label{kcs-so-do}

\rong{vtt-71-kcs-so-do.png}

Bấm mã lệnh để mở. Đầu trang là nhóm lệnh (lệnh chính và các lệnh đi kèm như ruột, tờ lót, gom chung một nhóm),
khách, sản phẩm và \oo{Đã kiểm đạt} trên số cần, cạnh đó là nút \oo{Đóng lệnh nhóm} (mục \ref{dong-lenh}).

Dải ô là các bước của lệnh. Mỗi ô ghi số tốt tổ đã làm, số lỗi đỏ nếu bước đó gây lỗi, và số tiêu chí bước cuối xét
cho công đoạn (hình danh sách kèm số). Ô ghi \q{lệnh chưa có tiêu chí} là công đoạn chưa khai tiêu chí, KCS vẫn ghi
lỗi ở dòng Lỗi khác. Ô thuê ngoài ghi \oo{Thuê ngoài}: KCS không kiểm tại chỗ bước thuê ngoài. Ô cuối là \oo{Nhập
kho} với số đã gửi, chưa gửi.

Bấm một ô, phần dưới đổi theo bước đó. Mở lệnh lần đầu thì phần dưới là bước cuối.

\Needspace*{33\baselineskip}
\subsection{Ghi lỗi ở bước giữa}\label{ghi-loi-giua}

\rong{vtt-72-kcs-cong-doan-giua.png}

Ở bước giữa, phần dưới có tổng Tổ đã làm, \oo{Lỗi do In gây}, \oo{Bắt ở} và ba tab: \oo{Tiêu chí} (lỗi ghi ở đây
và kết quả kiểm cuối của từng tiêu chí), \oo{Quy cách}, \oo{Lỗi đã ghi}. Thấy lỗi thì bấm \oo{Ghi lỗi}.

\rong{vtt-73-kcs-ghi-loi.png}

\so{1} Tick tiêu chí không đạt. Tiêu chí đạt để trống. Lỗi không có trong danh sách thì tick dòng \oo{Lỗi khác},
khi đó Ghi chú là bắt buộc.
\so{2} \oo{Số lỗi} theo đơn vị của bước.
\so{3} \oo{Ghi chú}: lỗi ở đâu, ví dụ góc phải tờ.
\so{4} \oo{Ảnh}: \oo{Chụp ảnh} (điện thoại mở thẳng camera) hoặc \oo{Chọn ảnh}. Ở bước giữa ảnh không bắt buộc.
Bấm \oo{Xong} để thu dòng lại.
\so{5} \oo{Lưu lỗi}. Chân hộp cộng Tổng lỗi lần này và báo lỗi sẽ về tổ nào.

Cột phải \oo{Đối chiếu} chỉ để xem: quy cách và tệp của lệnh, để so hàng với mẫu ngay tại chỗ. Lỗi lưu xong báo
ngay về tổ gây lỗi và trừ khỏi số tổ đó còn ghi được.

\Needspace*{17\baselineskip}
\subsection{Tổ xem lỗi KCS}\label{to-xem-loi}

\rong{vtt-74-to-loi-cho-xem.png}

Tổ gây lỗi nhận thông báo ngay, đầu màn Việc của tổ hiện chip \so{1} \oo{Lỗi KCS chờ xem} kèm số việc. Dòng việc có
chấm đỏ cạnh tên công đoạn và chip \q{KCS: 6 lỗi}. Bấm chip để lọc, bấm việc để mở.

\rong[0.85\linewidth]{vtt-75-to-tab-kcs.png}

Mở từ lỗi KCS thì ngăn vào thẳng \so{1} tab \oo{KCS}. \so{2} Khối \oo{Lỗi KCS}: tiêu chí, ghi chú, số lỗi, người ghi,
lúc ghi, ảnh (bấm con mắt để xem, mũi tên để tải). Không có nút \q{Đã xem}: mở tab là tính tổ đã xem, chip
\q{đã xem} hiện kèm giờ, bên KCS thấy cột Tổ đã xem đổi theo.

\Needspace*{35\baselineskip}
\subsection{Kiểm hàng ở bước cuối}\label{kiem-hang}

Bước cuối có hàng tổ đã làm thì nút \oo{Kiểm hàng} kèm số chưa kiểm sáng lên. Mỗi lần kiểm là kiểm hết số đang chờ.

\rong{vtt-76-kcs-kiem-hang.png}

Phần tick tiêu chí gộp tiêu chí của cả chuỗi công đoạn, chia theo công đoạn (ảnh này là ba tiêu chí của In). Tick
dòng nào thì dòng đó mở ra ô Số lỗi, Ghi chú, ảnh. Tick tiêu chí của công đoạn nào thì lỗi tính cho công đoạn đó.

\so{1} Dòng \oo{Lỗi khác}: ghi chú bắt buộc, kèm \so{2} \oo{Lỗi do công đoạn nào} để quy đúng tổ gây lỗi. Ảnh này
ghi 2 cái bung keo gáy, quy cho Dán bìa folder.

Ở bước cuối mỗi dòng lỗi phải có \oo{ít nhất 1} ảnh. Ô \oo{Ghi chú cho lần kiểm} không bắt buộc.

\so{3} Chân hộp: \oo{Lần này kiểm} bao nhiêu, \oo{Đạt} (số chưa kiểm trừ số lỗi), \oo{Lỗi}.
\so{4} \oo{Lưu kết quả kiểm}.

Đã lưu thì không sửa được. Kiểm sai phải nhờ người có quyền xử lý, nên đếm kỹ trước khi lưu. Cột phải
\oo{Đối chiếu} có thêm tab \oo{Tổ đã ghi}: các lần tổ ghi sản lượng, để biết hàng đang kiểm là của lần nào.

\Needspace*{12\baselineskip}
\rong{vtt-77-kcs-sau-kiem.png}

Lưu xong, bước cuối cập nhật Đạt 995, Lỗi 5, Chưa kiểm 0. Tab \oo{Lần kiểm} có dòng mới: lúc kiểm, người kiểm,
đạt, từng dòng lỗi với công đoạn gây lỗi, lỗi gì, ảnh và Tổ đã xem. Ô In và ô Dán bìa folder trên dải hiện số lỗi
đỏ. Hai tổ đó nhận lỗi như mục \ref{to-xem-loi}.

\rong{vtt-78-kcs-loi-theo-cong-doan.png}

Tab \oo{Lỗi theo công đoạn} gom các lần kiểm của bước cuối theo công đoạn gây lỗi, tổ, lỗi gì, số lỗi, số lần.

% ============================================================
\Needspace*{13\baselineskip}
\section{Gửi nhập kho và đóng lệnh}\label{nhap-kho-dong-lenh}

Kiểm xong có hàng đạt thì nút \oo{Gửi nhập kho} kèm số cái hiện cạnh Kiểm hàng.

\canh{0.45}{vtt-79-kcs-gui-nhap-kho.png}{%
Bấm \oo{Gửi nhập kho}, hộp hỏi lại số đề nghị kho nhập và lệnh. Bấm \oo{Gửi yêu cầu}.

Yêu cầu nhập sang màn Yêu cầu nhập xuất của kho, nguồn ghi \q{KCS gửi thành phẩm}, nhập về Kho thành phẩm. Thủ kho
lập phiếu, ghi sổ là xong phần kho.}

\rong{vtt-80-kcs-nhap-kho.png}

Tab \oo{Nhập kho} của bước cuối theo dõi từng yêu cầu: mã yêu cầu (bấm để mở bên màn Kho), gửi lúc, đề nghị, kho đã
nhận, trạng thái. Ảnh trên là lúc kho đã ghi sổ: \oo{Kho đã nhận đủ}. Chưa ai có ô Thao tác của KCS thì không ai gửi
được thành phẩm sang kho, đơn bán đứng mãi ở bước chờ KCS.

\Needspace*{16\baselineskip}
\subsection{Đóng lệnh}\label{dong-lenh}

\rong{vtt-81-kcs-dong-lenh.png}

Bấm \oo{Đóng lệnh nhóm} ở đầu trang. Khung xác nhận bật ngay dưới nút, nêu những điểm chưa xong (ảnh này: đạt 995
trên 1.000 cái, thiếu 5). Bấm \oo{Xác nhận đóng}, hoặc \oo{Để sau}.

Lệnh còn điểm chưa xong (thiếu số so với đơn, còn hàng chưa kiểm, chưa gửi kho, còn việc đang làm) thì chỉ người
có ô \oo{Đóng lệnh thiếu} đóng được. Người không có ô này thấy nút mờ kèm câu \q{Còn N điểm chưa xong, cần quyền
Đóng lệnh thiếu.}

\rong{vtt-82-kcs-da-dong.png}

Đóng xong, nhóm lệnh ghi \oo{Đã đóng} kèm người đóng và giờ. Nút đổi thành \oo{Mở lại nhóm}, cũng cần ô Đóng lệnh
thiếu. Mở lại thì việc chưa làm hiện lại ở Việc của tổ. Phía Kế hoạch sản xuất xem lệnh đã đóng ở
\kh{lenh-xong}{Khi lệnh xong}.

\Needspace*{27\baselineskip}
\subsection{Báo cáo lỗi}\label{bao-cao-loi}

Tab \oo{Báo cáo lỗi} của màn KCS tổng hợp lỗi theo kỳ, mặc định Tháng này theo ngày kiểm.

\rong{vtt-83-kcs-bao-cao-loi.png}

Đầu tab có ô tìm, kỳ, Lọc, ba cách gom \oo{Theo công đoạn}, \oo{Theo tổ}, \oo{Theo tiêu chí}, và hai kiểu xem
\oo{Bảng}, \oo{Biểu đồ}. Dải số: Lượt kiểm, Tỷ lệ đạt ở bước cuối, Lần ghi lỗi, Lệnh có lỗi, mỗi số so với kỳ
trước. Kiểu Bảng có bảng gom (mỗi ô là số lỗi trong ngày, đếm theo đơn vị nơi bắt, bấm một dòng để lọc) và bảng
\oo{Các lần ghi lỗi} mới nhất trước: lúc, lệnh, khách, công đoạn gây lỗi, tiêu chí, bắt ở đâu, số lỗi, tổ chịu,
ảnh, người ghi.

\rong{vtt-84-kcs-bieu-do.png}

Kiểu Biểu đồ: lần ghi lỗi mỗi ngày theo công đoạn gây lỗi (bấm một cột để xem riêng ngày đó), công đoạn gây lỗi
nhiều nhất, lỗi bắt tại công đoạn hay lọt tới KCS cuối, và tỷ lệ đạt ở bước cuối. Nút \oo{Xuất Excel} ở đầu màn
xuất đúng phần đang lọc.

% ============================================================
\Needspace*{10\baselineskip}
\section{Câu hỏi hay gặp}\label{hoi-dap}

\begin{description}[style=nextline,leftmargin=0pt,font=\normalfont\tieude\small]
\item[Bấm Bắt đầu không được?]
Đọc khung vàng \q{Chưa bắt đầu được vì} ở cột phải của ngăn. Bốn lý do và cách gỡ
ở mục \ref{chua-bat-dau}.

\item[Tổ trước nói đã giao mà không thấy hàng?]
Mở tab Giao và nhận. Dòng Nhận vào còn \q{chờ xác nhận} thì tổ
trưởng bấm Xác nhận. Không có dòng nào là tổ trước chưa ghi sản lượng hoặc chưa đề xuất giao, hỏi tổ trước. Hàng
thuê ngoài thì hỏi Kế hoạch sản xuất đã nhận về chưa.

\item[Ghi nhầm số?]
Tổ trưởng mở dòng Lần N ở tab Sản lượng, bấm Sửa số, ghi lý do (mục \ref{sua-so}).

\item[Không bấm được Kết thúc?]
Còn hàng tốt chưa đề xuất giao. Vào tab Giao và nhận bấm Đề xuất giao.

\item[Không thấy nút Trả khuôn về kệ?]
Nút chỉ hiện sau khi việc đã Kết thúc.

\item[Máy hư giữa ca?]
Báo dừng máy, tick \q{Máy đang dừng, không chạy được} để việc tự tạm dừng. Tổ còn máy khác
chạy được công đoạn này thì Đổi máy (chương \ref{dung-ket-thuc}).

\item[Thiếu giấy giữa chừng?]
Tab Vật tư, Yêu cầu bổ sung, gõ số xin thêm (chương \ref{vat-tu}).

\item[KCS ghi lỗi tổ mình mà tổ thấy không đúng?]
Mở tab KCS xem ảnh và ghi chú, rồi trao đổi với KCS. Số đã lưu
KCS không tự sửa được.

\item[Làm sao biết tổ đã xem lỗi chưa?]
Bên KCS, cột Tổ đã xem ở tab Lần kiểm và tab Lỗi đã ghi. Tổ mở tab KCS
của việc là tính đã xem.

\item[Lệnh thiếu vài cái mà khách nhận luôn?]
Người có ô Đóng lệnh thiếu bấm Đóng lệnh nhóm, khung xác nhận nêu
số thiếu, bấm Xác nhận đóng.

\item[Công đoạn nào cũng ghi \q{lệnh chưa có tiêu chí}?]
Danh mục Tiêu chí KCS chưa khai cho công đoạn đó. Nhờ
người khai danh mục thêm (\kb{danh-muc-kcs}{Tiêu chí KCS}). Lệnh đã phát hành lấy bộ mới bằng nút Cập nhật tiêu chí ở
hồ sơ lệnh (\kh{ho-so-lenh}{Hồ sơ lệnh sản xuất}).
\end{description}

\end{document}
